\documentclass[../../../../main.tex]{subfiles}
\begin{document}

\begin{flushright}
\footnotesize\textit{【自動文字起こし・要確認】}
\end{flushright}

\begin{tcolorbox}[title=整関数と割り切れに関するメモ]
\begin{itemize}
    \item $[ \text{整関数} ]$\\
    任意の $n$ 次整関数は
    \[
    f(x) = \sum_{k=0}^n f_k(x) a_k \quad \left( f_k(x) = (x+k)(x+k-1) \dots (x+1) x \right)
    \]
    と表せる.\\
    $f(x)$ が整関数 $\iff$ ある $k$ について $f(k) \in \mathbb{Z} \land f(x+1)-f(x)$ が整関数.
    \item $[ \text{割り切る} ]$\\
    $f(x)$ が $(x-\alpha)^k$ で割り切れるとき
    \[
    \iff f(\alpha) = f'(\alpha) = \dots = f^{(k-1)}(\alpha) = 0
    \]
    \item $[ \text{付} ]$\\
    $x-\alpha$ で割り切るための作業が出ている時, $t = x-\alpha$ とおいた方が見通しがよくなる.
\end{itemize}
\end{tcolorbox}

[解] $f_k(x) = x^k - 1 - k(x-1)$ とおくと 任意の $n$ 次多項式 $g(x)$ は,
\[
g(x) = \sum_{k=2}^n a_k f_k(x) + a_1 x + a_0
\]
とおける.

(i) $g(x)$ が $(x-1)^2$ で割り切れるには, $g(1)=g'(1)=0$ が必要十分であるが, $f_k(1)=f'_k(1)=0$ から, この条件は
\[
\begin{cases} a_1 + a_0 = 0 \\ a_1 = 0 \end{cases} \iff a_1 = a_0 = 0
\]
より, この時,
\[
g(x) = \sum_{k=2}^n a_k f_k(x) \quad \text{の形で表せる.} \quad \dots \text{①}
\]
逆に, ①の時, $g(1)=g'(1)=0$ となり, 十分.
以上から示された.

(ii) $g(x)$ が $(x-1)^3$ で割り切れるには, $g(1)=g'(1)=g''(1)=0$ が必要十分である. $f_k(1)=f'_k(1)=0, f''_k(1)=k(k-1)$ から
この時,
\[
\begin{cases} a_1 = a_0 = 0 \\ \sum_{k=2}^n k(k-1) a_k = 0 \end{cases}
\]
だから, たしかに $g(x) = \sum_{k=2}^n a_k f_k(x)$ ($\sum_{k=2}^n \frac{k(k-1)}{2} a_k = 0$) を書けることが必要. 逆に, このように書ける時, 同様にして十分である.
以上から示された.

[別解] (1) 十分性について.
\[
f_k(x) = \{(x-1)+1\}^k - k(x-1) - 1
\]
を $x-1$ の整式とみると
\[
f_k(x) = (t+1)^k - kt - 1 \quad (t=x-1)
\]
の 1 次以下の部分は $0$ だから, $f_k(x)$ は $(x-1)^2$ で割り切れる. よって, 十分.

・必要性
$n=1,2$ の時自明. $n < m$ の時, 「$(x-1)^2$ で割り切れる $g(x)$ は $\sum_{k=2}^n a_k f_k(x)$ の形で表せる」が真とする. $m+1$ 次の $g(x)$ は $m$ 次以下の整式 $h(x)$ として
\[
g(x) = a_m f_m(x) + h(x)
\]
と表せる. この時, $f_m(x), g(x)$ が共に $(x-1)^2$ で割り切れるので, $h(x)$ もそうである. よって仮定から示された.

(2) より, $f_k(x) = \sum_{i=2}^k {}_k\mathrm{C}_i (x-1)^i$ だから,
\[
\sum_{k=2}^n a_k f_k(x) = \sum_{k=2}^n a_k \left[ \sum_{i=2}^k {}_k\mathrm{C}_i (x-1)^i \right]
\]
で, これを $(x-1)^3$ でわったあまりは,
\[
\sum_{k=2}^n a_k {}_k\mathrm{C}_2 (x-1)^2
\]
だから, 示すべきことは明らか.

\end{document}
